\documentclass[12pt]{article} \usepackage[utf8x]{inputenc}
\usepackage{graphicx}
\usepackage{multirow}
\usepackage{hhline}
\usepackage{booktabs}
\usepackage{vmargin} %cambia el margen
\usepackage{amsmath, amsthm}
\usepackage{amsfonts}
\usepackage{float}
\usepackage{listings}
\usepackage{algorithm} % http://ctan.org/pkg/algorithms
\usepackage{algpseudocode} % http://ctan.org/pkg/algorithmicx


\usepackage[hidelinks]{hyperref}




\title{
  Historia de las Matemáticas\\
  \large Historia de los algoritmos inspirados en la evolución:
  Algoritmos Genéticos, Programación Evolutiva, Estrategias de
  Evolución y Programación Genética.  }


\author{
  Francisco Luque Sánchez \\
  Ignacio Mas Mesa \\
  Miguel Morales Castillo \\
  María del Mar Ruiz Martín \\
}



\begin{document}
\maketitle
\begin{center}
\includegraphics[scale=0.35]{escudo.png}
\end{center}

\newpage

\tableofcontents % para generar el índice de contenidos

\pagebreak

\section{Introducción}

Las metaheurísticas son una clase de estrategias de resolución de
problemas que aparecieron en los años 50. Este tipo de algoritmos se
inspiran en procesos naturales, físicos o sociales para intentar
resolver problemas, usualmente de optimización, cuando el espacio de
soluciones es demasiado amplio para ser estudiado por métodos
tradicionales.\\

En este trabajo abordaremos, desde una perspectiva histórica, el
origen, evolución y estado actual de este tipo de técnicas, así como
ejemplos prácticos de aplicación. En particular, incidiremos sobre una
clase importante de metaheurísticas, las cuales se inspiran en el
proceso natural de la evolución. Dentro de este grupo se engloban
cuatro enfoques principales, que son los algoritmos genéticos, la
programación evolutiva, la programación genética y las estatregias
de evolución.\\

\section{Metaheurísticas}

Las metaheurísticas aparecieron en los años 50 como respuesta a los
problemas de optimización. Hasta ese momento, los ordenadores habían
sido utilizados como herramientas de ayuda al cálculo, pero no se
habían pensado como máquinas capaces de resolver problemas más
complejos (needs further research). Cuando se intentaron implementar
técnicas de optimización clásica para resolver problemas de búsqueda
(derivación de funciones para el cálculo de extremos locales), se
observó que para muchos problemas este enfoque no era suficiente
(citation needed). Uno de los problemas que surgió fue que las
funciones a optimizar tenían muchos extremos locales, infinitos en
algunos casos. También se encontraron problemas en los que las
funciones a optimizar no eran derivables, por lo que el enfoque
clásico no podía
ser aplicado.\\

Empezaron a aparecer entonces intentos de resolución de problemas de
este tipo con estrategias distintas. El primer ejemplo de este tipo de
estrategias vino en 1951, con un método de aproximación estocástica.
Este método, creado por Herbert Robins y Sutton Monro, trataba de
encontrar extremos de una función utilizando un método de aproximación
por pasos que depende de una variable aleatoria. Así se puede buscar
el mínimo de una función sin tener que buscar la derivada de la misma.\\

\section{Algoritmos genéticos}

\section{Programación evolutiva}

\section{Estrategias de evolución}

De forma casi paralela a la aparición de los \textbf{algoritmos
 genéticos} y la \textbf{programación evolutiva} en Estados Unidos,
se desarrollaron en la segunda mitad de la década de los 60 en
Alemania las \textbf{estrategias evolutivas} (\textit{evolution
  strategies}, ES). En 1964, tres estudiantes de la Technical University of
Berlin, viz.: Bienert, Rechenberg y Schwefel, se encontraron con un
problema de diseño durante sus estudios de tecnología aeroespacial. La
pieza a diseñar debía ser probada en una serie de experimentos en el
túnel de viento con el objetivo de minimizar su rozamiento. \\

Tras aplicar diversas técnicas \textit{clásicas} de análisis numérico a
estos efectos, como el gradiente descendente, los resultados obtenidos
fueron pobres, al menos lo suficiente como para tratar de explorar
otras posibilidades. En ese momento Rechenberg decidió aplicar
técnicas estocásticas usando dados para la generación de números
aleatorios. Así nació la primera versión de las estrategias
evolutivas, una $(1 + 1)$ ES en la notación actual, i.e.: en cada
paso se consideraba una única solución \textit{padre}, de la cual se
variaban aleatoriamente ciertas características. Una vez hecho esto,
se comparaba la versión modificada con la original y se conservaba la
mejor de ambas, continuando con el proceso mientras se considerara
conveniente. \\

Para comprobar la eficacia de esta nueva aproximación se trató de
encontrar el óptimo, en un problema en dos dimensiones, de un problema
similar al que se enfrentaban. Durante unas cuantas iteraciones el
proto-algoritmo se mantuvo encallado en un óptimo local, para
finalmente converger al óptimo global teórico. \\

Más tarde, Bienert construyó un robot para llevar a cabo estos
cálculos de forma automática; y Lihchtfu{\ss} aplicó con éxito esta técnica para
el diseño de una tubería. El resultado fue inesperado, pero claramente
mejor que los diseños conocidos hasta el momento. Asimismo, se empleó
para diseñar una tobera que debía propulsar agua a velocidades
supersónicas. Este experimento resultó especialmente interesante pues,
entre otros factores, se desconocía la longitud óptima de la
tobera. Para incluir esto en el espacio de búsqueda se emuló el
duplicado/eliminación de genes, de forma que se pudiera trabajar con un número
variable de parámetros indicando el diámetro de la tobera en distintos
puntos. \\

Sin embargo, pronto se descubrió que esta estrategia tenía sus
inconvenientes. En particular, las distribuciones binomiales usadas
hasta el momento producían, en muchas ocasiones, convergencia
prematura a óptimos locales que dejaban mucho que desear. Es por esto
que de ahí en adelante se adoptó como práctica común el empleo de
distribuciones normales en lugar de binomiales para la generación de
números aleatorios. \\

Posteriormente, Rechenberg elaboró su tesis en 1971 estudiando los
ratios de convergencia de las soluciones en distintos casos,
obteniendo la famosa \textit{regla del quinto} (\textit{$\frac 1 5$th
  rule}) y trató de optimizar el algoritmo aproximándolo más a lo que
ocurría en la naturaleza, considerando una \textit{población} de
soluciones en cada paso, introduciendo el paso de la
\textit{recombinación} de soluciones además de la mutación y dando
lugar a lo que, en la notación de Schwefel se denominó $(\mu + 1)$ ES
($\mu$ padres, $1$ hijo). Nótese que esto no habría sido posible en
las $(1 + 1)$ ES primigenias. \\

No obstante, aun esta modificación no fue suficiente en términos de
capacidad de adaptación, por lo que aparecieron las $(\mu + \lambda)$
y $(\mu, \lambda)$ ES, con $\mu$ padres y $\lambda$ descendientes en
cada generación. \\

\begin{description}
\item[$(\mu + \lambda)$] Se seleccionan los $\mu$
mejores individuos de entre los $\mu$ padres y los $\lambda$
descendientes de cada generación.
\item[$(\mu, \lambda)$] También se crean $\lambda$ descendientes en
  cada generación a partir de los $\mu$ padres. En este caso, sin
  embargo, los padres \textbf{no} pueden sobrevivir, seleccionándose
  los $\mu$ mejores descendientes de los $\lambda$
  considerados. Claramente, es necesario que $\lambda > \mu$. En el
  caso límite $\mu = \lambda$ todos los descendientes de cada
  generación se seleccionan como padres, sin importar su bondad, por
  lo que la población simplemente se dispersa aleatoriamente por el
  espacio de búsqueda.
\end{description}

Cabe destacar que, pese a los éxitos cosechados experimentalmente, los
inicios de esta disciplina fueron duros pues la comunidad científica
de la época se mostraba reticente a aceptar (creer) su versatilidad y
eficiencia por varios motivos.

\begin{itemize}
\item Durante la década de los 60 los métodos numéricos observaron un
  desarrollo inaudito, ganándose el respeto de los investigadores de
  la época y convirtiéndose incluso en \textit{buzzword}. Así, la
  necesidad de considerar otras técnicas era cuestionable,
  especialmente estas que no venían acompañadas de una justificación
  teórica en cuanto a su eficacia. Esto cambió más tarde tras diversos
  análisis por parte de Rechenberg, Schwefel y otros sobre su
  algoritmo.
\item Aun aceptando la posibilidad de que la $(1 + 1)$ ES fuera un
  artefacto útil en ciertos casos, la incorporación de las $(\lambda +
  \mu)$ ES no tuvo muy buena acogida. Por un lado, considerar $\mu >
  1$ individuos en la población se consideraba innecesario, pues
  simplemente aumentaba los tiempos de cálculo arrastrando individuos
  poco (menos) prometedores. Por otro, $\lambda > 1$ parecía no
  explotar inmediatamente el conocimiento obtenido en generaciones previas.
\end{itemize}

Si bien la crítica referida a los $\lambda > 1$ padres podía ser
rebatida gracias a la computación en paralelo de los descendientes, el
contraargumento no tenía mucho peso pues los ordenadores con
arquitecturas capaces de procesamiento en paralelo no existían ni se
los esperaba en el futuro próximo. \\

Finalmente, tras todos los cambios introducidos a lo largo de su
historia, la versión actual de una ES genérica puede verse resumida en la
siguiente figura:

\begin{figure}
\begin{algorithm}[H]
  \begin{algorithmic}[1]
    \label{ES-alg}
    \Procedure ES
    \State $g := 0$
    \State initialize $\left( P_p^{(0)} := \left\{ \left(y_m^{(0)},
            s_m^{(0)}, F(y_m^{(0)}) \right), m = 1, ..., \mu \right\}
    \right);$
    \Repeat

    \For {$l := 1$ To $\lambda$}
    \State $O_l :=$ marriage $\left( P_p^{(g)}, \rho \right);$
    \State $s_l := s_{recombination} \left( O_l \right);$
    \State $y_l := y_{recombination} \left( O_l \right);$
    \State $\tilde{s_l} := s_{mutation} \left( s_l \right);$
    \State $\tilde{y_l} := y_{mutation} \left( y_l, \tilde{s_l} \right);$
    \State $\tilde{F}_l := F(\tilde{y_l})$
    \EndFor

    \State $P_0^{(g)} := \left\{ \left( \tilde{y_l}, \tilde{s_l},
          \tilde{F_l} \right), l = 1, ..., \lambda \right\};$

    \If {$selection_{type} == (\mu, \lambda)$}
    \State $P_p^{(g+1)} := selection \left( P_o^{(g)}, \mu \right);$
    \ElsIf {$selection_{type} == (\mu + \lambda)$}
    \State $P_p^{(g+1)} := selection \left( P_o^{(g)}, P_p^{(g)} \mu \right);$
    \EndIf

    \State $g := g + 1$

    \Until termination condition
    \EndProcedure
  \end{algorithmic}
\end{algorithm}
  \caption{Pseudocódigo del $(\mu / \rho, {+}\over{,} \lambda)$-ES \cite{paper-es-salva}}
\end{figure}

En la notación actual, las $(\mu / \rho {+}\over{,} \lambda)$ ES consideran $\mu$ padres en cada generación, escogiéndose $\lambda$ subconjuntos de $\rho$ individuos para dar, cada uno de ellos, uno de los $\lambda$ descendientes. Como se había discutido previamente, en la versión ``$+$'' tanto los padres como la prole pueden pasar a la siguiente generación, mientras que en la versión ``,'' sólo los descendientes pueden sobrevivir. \\

La descripción del pseudocódigo del algoritmo es como sigue:

\begin{enumerate}
\item En la primera generación ($g = 0$) se inicializa la población de padres (de qué forma concreta dependerá de cada problema).
\item A continuación entramos en el bucle principal
  \begin{enumerate}
  \item Se comienza generando los descendientes de la generación actual. Cada iteración genera un nuevo vástago.
    \begin{enumerate}
    \item En primer lugar, el paso de la maridación (\textit{marriage}) selecciona una \textit{familia} de padres $O_l$ de tamaño $\rho$. Esta selección es completamente aleatoria, a diferencia de los algoritmos genéticos en los que se favorece la selección de los mejores padres.
    \item Más tarde le llega el turno a la recombinación, tanto de las características de las soluciones (\textit{object parameters}) como de las de la estrategia (\textit{strategy parameters}). De nuevo, cómo se recombine la información de cada padre dependerá de una implementación a otra y no existe una forma óptima.
    \item Por último, se mutan los \textit{strategy parameters} obtenidos de la recombinación y, sólo entonces, las características de la solución con los nuevos \textit{strategy parameters} obtenidos. Finalmente se calcula la bondad de la nueva solución obtenida mediante la función objetivo.
    \end{enumerate}
  \item Una vez generados todos los descendientes, se distingue entre el esquema de  selección $(\mu, \lambda)$ y $(\mu + \lambda)$:
    \begin{enumerate}
    \item En el primer caso se seleccionan los $\mu$ descendientes
    \item En el segundo caso se seleccionan los $\mu$ mejores individuos de entre los padres ($P_p^{(g)}$) y los descendientes ($P_o^{(g)}$).
    \end{enumerate}
  \item Para finalizar se actualiza la generación actual.
  \item El bucle principal se repetirá mientras no se alcance la condición de parada que, de nuevo, depende en gran medida de la casuística concreta del problema.
  \end{enumerate}
\end{enumerate}

\section{Programación genética}

\section{Referencias}

\end{document}
